% ============================================================================
% Fichier  : partie6/P06_C17_profilage_ia.tex
% Titre    : Profilage, Décisions Automatisées et Intelligence Artificielle
% Auteur   : MINKA MI NGUIDJOÏ Thierry Emmanuel
% Version  : 1.0 -- novembre 2026
% Corrigés : corriges/Corr_C17.tex
% ============================================================================
% Sources : loi n° 2024/017, art. 5 (profilage), 9(3), 17, 18, 21, 44, 48,
% 65 (texte officiel) ; ISO/IEC 27701 ; GCI C19 (gouvernance de l'IA).
% Faits nouveaux (bible) : NOUN SCORE ; refus de crédit de M. ATANGANA du
% 12/08/2026 ; selfies de vérification d'identité ; pointage biométrique de
% MBOLO ; audit de NOUN PAY, octobre 2026.
% ============================================================================

\chapter{Profilage, Décisions Automatisées et Intelligence Artificielle}
\label{chap:profilage}
\chaptermark{Profilage et IA}

\epigraph{%
    « Un modèle qui décide d'un crédit\\
    doit pouvoir être expliqué à celui qu'il refuse. »%
}{Principe du cours}

\niveauChapitre{Expert}{%
    Chapitres~\ref{chap:principes}, \ref{chap:risque}, \ref{chap:droits}
    et~\ref{chap:transferts}
}

\begin{boxobjectifs}
\begin{itemize}
    \item Reconnaître un profilage et une décision fondée exclusivement sur
          un traitement automatisé.
    \item Appliquer l'article~44 et vérifier les garanties dues à la
          personne~: intervention humaine, point de vue, contestation.
    \item Lire l'article~65 et dire ce que l'auditeur peut en conclure.
    \item Auditer un modèle de notation~: données d'entrée, finalité,
          information, non-discrimination.
    \item Traiter les données biométriques au regard des articles~5 et~48.
\end{itemize}
\end{boxobjectifs}

\bigskip

% ============================================================================
\section*{Avant de commencer : « l'ordinateur a dit non »}
% ============================================================================

Bafoussam, mercredi 12~août~2026. M.~ATANGANA, couturier, demande à NOUN
PAY un microcrédit de 150\,000~francs CFA pour acheter une machine. La
réponse arrive en quatre secondes~: refusé. Il appelle le service client,
qui lui explique que « l'ordinateur a dit non » et qu'aucun agent ne peut
revenir sur la décision. Il écrit pour demander pourquoi.

Personne, à NOUN PAY, ne sait le lui dire. Le modèle NOUN SCORE a été
entraîné en 2025 par un prestataire sur les historiques de 210\,000
utilisateurs. Il utilise une soixantaine de variables~: régularité des
recharges de crédit téléphonique, montants reçus, heure habituelle des
transactions, nombre de contacts dans le répertoire du téléphone, et
localisation approximative. L'équipe produit connaît les variables~; elle
ne sait pas lesquelles ont pesé pour M.~ATANGANA.

\begin{boxalert}
\textbf{Violation typique.} Laisser un modèle décider seul d'une question
importante pour une personne, sans pouvoir expliquer la décision ni la
faire réexaminer par un humain.

\smallskip
\textbf{Règle~:} la personne a le droit de s'opposer à une décision fondée
exclusivement sur un traitement automatisé, et, en tout cas, d'obtenir une
intervention humaine, d'exprimer son point de vue et de contester
(article~44). Elle doit avoir été informée de la logique sous-jacente
(article~21).
\end{boxalert}

% ============================================================================
\section{Profilage et Décision Automatisée}
\label{sec:c14-definitions}
% ============================================================================

\begin{boxdef}
\termecle{Profilage}~: selon l'article~5, « traitement automatisé de données
à caractère personnel consistant à les utiliser pour évaluer certains
aspects personnels relatifs à une personne physique, notamment sa santé,
ses préférences, sa localisation et sa situation économique ».
\end{boxdef}

Le profilage \textit{évalue}~; la décision automatisée \textit{décide}. Un
score de solvabilité est un profilage~; le refus de crédit prononcé par le
système sans intervention humaine est une décision automatisée. Les deux se
recouvrent souvent, mais pas toujours~: un score présenté à un agent qui
décide réellement est un profilage sans décision automatisée.

\begin{boxjuris}
\textbf{\loicam{n°~2024/017}, article~44}~\cite{loi2024017}~:
\textit{« (1) La personne concernée a le droit de s'opposer à toute décision
basée exclusivement sur le traitement automatisé de ses données à
caractère personnel, y compris le profilage. (2) Les dispositions de
l'alinéa 1 ci-dessus ne s'appliquent pas si~: la personne concernée est
informée de l'utilisation du système automatisé de prise de décision et a
donné son consentement préalable, explicite et éclairé~; le traitement est
autorisé par la loi, sous réserve qu'elle prévoit des mesures appropriées
pour la sauvegarde des droits et libertés et des intérêts légitimes de la
personne concernée. (3) La personne concernée a également le droit~:
d'obtenir une intervention humaine de la part du responsable du
traitement~; d'exprimer son point de vue~; de contester la décision fondée
sur le traitement automatisé. »}

\smallskip
\textbf{Portée.} L'alinéa~3 ne connaît pas d'exception~: même quand la
personne a consenti, elle conserve le droit à une intervention humaine et
à la contestation. Un système « non révisable » méconnaît l'article~44.
\end{boxjuris}

% ============================================================================
\section{L'Article 65}
\label{sec:c14-article65}
% ============================================================================

\begin{boxjuris}
\textbf{\loicam{n°~2024/017}, article~65}~\cite{loi2024017}~:
\textit{« Est puni d'un emprisonnement de trois (03) à dix (10) ans et
d'une amende de un million (1~000~000) à vingt millions (20~000~000) de
francs CFA, ou de l'une de ces deux peines seulement, le responsable de
traitement ou le sous-traitant qui procède ou fait procéder à un
traitement de données à caractère personnel à des fins de profilage. »}
\end{boxjuris}

Pris isolément, le texte punit tout profilage. Lu avec l'article~44, qui
admet les décisions automatisées, y compris le profilage, sous conditions,
il pourrait viser seulement le profilage réalisé hors de ces conditions.
Les deux lectures sont défendables~; seul un juge, ou un texte
d'application, les départagera.

L'auditeur ne tranche pas. Il fait trois choses. Il signale la difficulté
au commanditaire, en citant les deux articles. Il cote \NEK{} de critère la
proposition « le profilage pratiqué est licite au regard de l'article~65 ».
Et il vérifie, sur pièces, toutes les conditions qui rendraient le
profilage défendable dans la lecture la plus favorable~: autorisation
demandée, consentement préalable, explicite et éclairé, information sur la
logique, intervention humaine effective, absence de discrimination. Si l'une
manque, la recommandation ne dépend plus de la lecture de l'article~65.

% ============================================================================
\section{Auditer un Modèle de Notation}
\label{sec:c14-modele}
% ============================================================================

L'audit d'un modèle ne demande pas d'être statisticien~; il demande de
poser, sur chaque élément du modèle, les questions de la loi.

\begin{tableaucours}{L{3.2cm}L{4.8cm}L{4.0cm}}{%
    \tableauHeader{Élément} &
    \tableauHeader{Question d'audit} &
    \tableauHeader{Articles}
}
Données d'entraînement & Collectées pour cette finalité, ou détournées~?
Consentement~? & 9, 11, 67 \\
Variables & Chacune a-t-elle un lien avec la solvabilité~? Une variable
révèle-t-elle une donnée sensible ou une origine~? & 11, 17, 48 \\
Décision & Exclusivement automatisée~? Intervention humaine possible et
effective~? & 44 \\
Information & La personne est-elle informée de la logique et des
conséquences~? & 21 \\
Non-discrimination & Les taux de refus diffèrent-ils selon la région, la
langue, le sexe~? & 17 \\
Prestataire & Le modèle est-il hébergé ou réentraîné à l'étranger~? & 30,
32 \\
\end{tableaucours}
\vspace{-0.8em}
\begin{center}{\small\itshape Tableau~17.1. Les questions de la loi sur un modèle de notation}\end{center}

\subsection{Les variables qui révèlent}

Une variable anodine peut révéler une donnée sensible. La localisation
habituelle peut révéler l'origine régionale ou le lieu de culte~; l'heure
des transactions, la religion pendant un jeûne~; le répertoire des
contacts, des données sur des tiers qui n'ont jamais consenti. L'article~17
impose de traiter les données « sans tenir compte de la classe sociale,
de l'origine ethnique ou régionale, de l'appartenance syndicale, de
l'opinion politique ou de la conviction religieuse de la personne
concernée »~: une variable qui sert
d'intermédiaire à l'une d'elles fait entrer ce critère dans la décision.

\subsection{Mesurer une discrimination}

L'auditeur compare les taux de refus entre groupes de personnes. Si les
demandeurs d'une région sont refusés deux fois plus que ceux d'une autre à
revenus comparables, le modèle traite probablement une origine régionale
par intermédiaire. Cette comparaison exige de connaître le groupe des
personnes~; si l'organisation ne le connaît pas, et elle ne doit pas
collecter une donnée sensible pour le savoir, l'auditeur utilise un
indicateur indirect déclaré (la région de l'agence d'inscription) et cote
avec prudence. Le laboratoire APDC-LAB-10 conduit ce calcul.

% ============================================================================
\section{Les Données Biométriques}
\label{sec:c14-biometrie}
% ============================================================================

La biométrie figure dans la liste des données sensibles de l'article~5, et
l'article~48, alinéa~1, interdit le traitement des données relatives
« à la génétique et à la biométrie en matière de santé ». Le
Chapitre~\ref{chap:donnees} a relevé l'ambiguïté de cette fin de phrase.
Deux usages du cours la mettent à l'épreuve~:

\begin{itemize}
    \item le \textbf{selfie de vérification d'identité} de NOUN PAY, comparé
          automatiquement à la photographie de la pièce d'identité~;
    \item le \textbf{pointage par empreinte digitale} du personnel de MBOLO,
          en place depuis 2023, qui n'a rien de médical mais se déroule dans
          un établissement de santé.
\end{itemize}

Dans les deux cas, l'auditeur établit les faits (quelles données, où elles
sont stockées, combien de temps, qui y accède), vérifie le consentement
exprès que l'article~66 exige pour conserver des données sensibles, et cote
\NEK{} de critère la licéité au regard de l'article~48 tant qu'aucune
autorisation ni aucun texte ne l'éclaire. Il recommande dans tous les cas
une solution de rechange non biométrique et le stockage du gabarit sur un
support détenu par la personne plutôt que dans une base centrale.

% ============================================================================
\section{L'Intelligence Artificielle : Trois Textes à Connaître}
\label{sec:c17-ia}
% ============================================================================

La loi n°~2024/017 ne contient pas le mot « intelligence artificielle ».
Elle saisit ces systèmes par leurs données et par leurs décisions~: les
articles~44 et~65 pour la décision automatisée, les principes pour les
données d'entraînement. Trois textes extérieurs complètent ce regard.
Aucun n'est un critère obligatoire au Cameroun. L'auditeur doit pourtant
les connaître, parce que les organisations qu'il audite achètent des
modèles conçus sous leur empire.

\begin{tableaucours}{L{3.0cm}L{2.4cm}L{6.6cm}}{%
    \tableauHeader{Texte} &
    \tableauHeader{Nature} &
    \tableauHeader{Ce que l'auditeur en retient}
}
NIST AI RMF~1.0 (2023)~\cite{nist2023airmf} & Cadre volontaire & Quatre
fonctions~: gouverner, cartographier, mesurer, gérer. Une liste de
qualités attendues d'un système, dont l'équité, l'explicabilité et le
respect de la vie privée \\
Règlement (UE) 2024/1689~\cite{ue2024ia} & Texte obligatoire dans l'Union &
Une approche par niveau de risque~: pratiques interdites, systèmes à haut
risque soumis à des obligations, obligations de transparence. La notation
de solvabilité des personnes physiques figure parmi les usages à haut
risque \\
Stratégie continentale de l'Union africaine (2024)~\cite{ua2024ia} &
Orientation politique & Un appel aux États à se doter de cadres nationaux
et à protéger les données dans le développement de l'intelligence
artificielle \\
\end{tableaucours}
\vspace{-0.8em}
\begin{center}{\small\itshape Tableau~17.2. Trois textes sur l'intelligence artificielle}\end{center}

\subsection{Ce que ces textes changent pour l'audit}

Le règlement européen s'applique à qui met un système sur le marché de
l'Union ou l'y utilise. NOUN SCORE note des clients camerounais et n'en
relève pas. La question se poserait si le modèle servait à NOUN
DIASPORA~: l'équipe juridique devrait alors l'instruire, comme elle le
fait pour le RGPD (Chapitre~\ref{chap:rgpd}).

Le cadre du NIST fournit un vocabulaire et une liste de vérification. Il
aide à ne rien oublier~: qui répond du modèle, sur quelles données il a
appris, comment sa performance et ses écarts entre groupes sont mesurés,
qui décide de le retirer. L'auditeur peut l'emprunter comme critère
volontaire, à condition de le déclarer (Chapitre~\ref{chap:normes}).

La stratégie africaine ne crée aucune obligation pour une entreprise.
Elle annonce la direction que prendront les textes nationaux.

\begin{boxalert}
\textbf{Violation typique.} Écrire dans un rapport qu'un modèle « n'est
pas conforme au règlement européen sur l'IA » alors que ce règlement ne
s'applique pas à l'organisation.

\smallskip
\textbf{Règle~:} un texte étranger se cite pour ce qu'il est. S'il
s'applique, c'est un critère obligatoire et le rapport dit pourquoi.
S'il ne s'applique pas, c'est une référence empruntée, déclarée au plan,
qui ne fonde aucune non-conformité légale.
\end{boxalert}

\subsection{Les questions propres aux systèmes qui apprennent}

Un modèle qui apprend pose quatre questions que la notation classique ne
posait pas.

\begin{enumerate}
    \item \textbf{Les données d'entraînement.} D'où viennent-elles, et les
          personnes ont-elles été informées de cet usage (articles~9, 11
          et~21)~?
    \item \textbf{La dérive.} Le modèle est-il réentraîné, et qui vérifie
          que ses résultats ne changent pas en silence~?
    \item \textbf{L'explication.} L'organisation sait-elle dire à une
          personne pourquoi elle a été refusée, dans des mots qu'elle
          comprend (article~44)~?
    \item \textbf{Le fournisseur.} Si le modèle vient d'un tiers, que dit
          le contrat sur la réutilisation des données des clients pour
          améliorer le produit (articles~30 et~32)~?
\end{enumerate}

Quand l'organisation ne peut pas répondre, la cote est \NEK{} de preuve.
Elle se lève par des pièces~: la fiche du modèle, le journal des
réentraînements, le contrat.

% ============================================================================
\section{L'Audit de NOUN SCORE}
\label{sec:c14-nounpay}
% ============================================================================

\begin{boxscenario}{Mission NOUN PAY. La notation automatique, octobre 2026}
\textbf{Constatations.}
\begin{itemize}
    \item 100~\% des demandes de microcrédit sont décidées par NOUN SCORE
          sans intervention humaine~; aucune procédure de réexamen.
    \item Les conditions d'utilisation mentionnent « l'utilisation
          d'algorithmes » sans décrire la logique ni les conséquences~; le
          consentement est une case unique couvrant tous les services.
    \item Le modèle utilise le nombre de contacts du répertoire, collecté
          par l'application sans consentement distinct.
    \item Sur les 41\,200 demandes de juillet à septembre~2026, le taux de
          refus est de 38~\% pour les inscrits des agences de l'Ouest et de
          61~\% pour ceux de l'Extrême-Nord, à tranches de revenus déclarés
          comparables.
    \item Aucune demande d'autorisation, aucune évaluation d'impact.
\end{itemize}

\textbf{Analyse} (confiance moyenne sur la discrimination, élevée sur le
reste). L'écart de taux de refus, à revenus comparables, suggère que des
variables de localisation servent d'intermédiaire à une origine
régionale~; l'indicateur indirect retenu (l'agence d'inscription) ne permet
pas de l'établir avec certitude. Les autres écarts sont établis.
\end{boxscenario}

\begin{boxlecture}{NOUN SCORE, octobre 2026}
\begin{tableaucours}{L{7.8cm}cL{3.2cm}}{%
    \tableauHeader{Exigence} &
    \tableauHeader{Cote} &
    \tableauHeader{Preuve}
}
Une intervention humaine est possible (art.~44(3)) & \NON{} & Procédure
absente, refus à M.~ATANGANA \\
La personne est informée de la logique (art.~21) & \NON{} & Conditions
d'utilisation \\
Le consentement à la décision automatisée est explicite (art.~44(2)) &
\NON{} & Case unique \\
Les variables sont liées à la finalité (art.~11) & \NON{} & Répertoire des
contacts \\
Le modèle ne discrimine pas selon l'origine (art.~17) & \NEK{} & Écart de
refus, indicateur indirect \\
Une évaluation d'impact a été conduite (art.~33) & \NON{} & Aucune \\
Le traitement est autorisé (art.~19) & \NON{} & Aucune autorisation \\
Le profilage est licite au regard de l'article~65 & \NEK{} & Lecture non
tranchée \\
\end{tableaucours}

\medskip
\textbf{Décompte.} Zéro \OUI{}, six \NON{}, deux \NEK{} sur huit.
\hfill\lecture{0}{0{,}25}{0{,}81}

\medskip
\textbf{Lecture.} Aucune exigence n'est satisfaite. Les deux \NEK{} portent
sur les questions les plus graves, discrimination et délit, mais la
recommandation n'en dépend pas~: les six \NON{} suffisent à demander la
suspension des refus automatiques jusqu'à la mise en place d'une
intervention humaine.
\end{boxlecture}

% ============================================================================
\section{Ce que le Profilage demande à chaque Métier}
\label{sec:c14-metiers}
% ============================================================================

\begin{boxmetier}{Responsable de traitement}
Ne jamais mettre en production un modèle qui décide sans prévoir le
réexamen humain, l'explication et la contestation. M.~NJOYA a suspendu les
refus automatiques le 20~octobre~: tout refus est désormais revu par un
agent.
\end{boxmetier}

\begin{boxmetier}{Auditeur de la protection des données}
Poser les questions de la loi variable par variable, mesurer les écarts de
traitement entre groupes avec un indicateur déclaré, et signaler la
difficulté de l'article~65 sans la trancher.
\end{boxmetier}

\begin{boxmetier}{Juriste conformité}
Rédiger l'information sur la logique de la décision dans des termes que la
personne comprend, et prévoir le circuit de contestation.
\end{boxmetier}

\begin{boxresume}
\begin{itemize}
    \item Le profilage évalue~; la décision automatisée décide~; les deux
          se recouvrent souvent sans se confondre.
    \item L'article~44 garantit, même après consentement, l'intervention
          humaine, l'expression du point de vue et la contestation.
    \item L'article~65 punit le traitement à des fins de profilage~; sa
          portée exacte n'est pas tranchée, et l'auditeur le dit.
    \item Une variable anodine peut servir d'intermédiaire à une donnée
          sensible~; l'audit compare les taux de décision entre groupes.
    \item À NOUN PAY, aucune exigence n'est satisfaite~; les refus
          automatiques ont été suspendus.
\end{itemize}
\end{boxresume}

\begin{boxinfo}
\textbf{Points de contrôle.} Vous avez compris ce chapitre si vous savez~:
\begin{itemize}
    \item distinguer profilage et décision automatisée~;
    \item énoncer les droits de l'article~44(3)~;
    \item repérer une variable qui sert d'intermédiaire à une origine.
\end{itemize}
\end{boxinfo}

\begin{boxexo}[title={\ding{45}\quad Exercice~17.1. Profilage, décision, ni l'un ni l'autre \niveaubase}]
(a) Un score de solvabilité montré à un conseiller qui décide~; (b) un
refus de crédit automatique~; (c) une liste de clients triés par ancienneté~;
(d) une publicité ciblée selon la localisation~; (e) le blocage automatique
d'une transaction suspecte.
\end{boxexo}

\begin{boxexo}[title={\ding{45}\quad Exercice~17.2. Répondre à M. ATANGANA \niveaumoyen}]
Rédigez la réponse de NOUN PAY à M.~ATANGANA~: information sur la logique,
réexamen humain, possibilité de contester. Que lui dites-vous des variables
utilisées~?
\end{boxexo}

\begin{boxexo}[title={\ding{45}\quad Exercice~17.3. Les taux de refus \niveaumoyen}]
Calculez le rapport des taux de refus entre l'Extrême-Nord et l'Ouest. Que
faudrait-il connaître de plus pour passer la proposition « le modèle
discrimine selon l'origine » de \NEK{} à \OUI{} ou à \NON{}~? Quelles
données ne devez-vous pas collecter pour le savoir~?
\end{boxexo}

\begin{boxexo}[title={\ding{45}\quad Exercice~17.4. Le pointage biométrique \niveauexpert}]
Le Dr~ABENA veut conserver le pointage par empreinte « parce que les
salariés trichaient avec les badges ». Analysez le traitement au regard des
articles~5, 9, 48 et~66. Proposez deux solutions, et dites comment vous
coterez la licéité de chacune.
\end{boxexo}

\bigskip

\begin{boxinfo}
\textbf{Pour aller plus loin.}
\begin{itemize}
    \item L'ISO/IEC~27701~\cite{iso27701}, pour l'encadrement des
          décisions automatisées.
    \item Le cours de gouvernance des SI de l'auteur (UE MSN-2S3-MSI03),
          chapitre sur la gouvernance des données et de l'IA.
\end{itemize}
\end{boxinfo}

\begin{boxinfo}
\textbf{Transition.} La santé, la banque, la justice et l'état civil sont
renvoyés par la loi à des textes particuliers. Comment auditer dans ces
secteurs, quand le régime général et le régime spécial se rencontrent~?
C'est l'objet du Chapitre~\ref{chap:secteurs}.
\end{boxinfo}
